\section{Discussion}\label{sec:discussion}

\subsection{What the two-land model shows}

The two-land model yields three results relevant to human--environment models \citep{Schaefer1954,May1973,Meadows1972,Meadows2004,WackernagelRees1996,BranderTaylor1998}. First, conversion can add to total biocapacity while $A_c$ falls. Whether $B$ rises overall also depends on other land flows and quality change; Proposition~\ref{prop:mask} gives the exact condition for the conversion term. The technology boost is capped, while debt depends on the magnitude and duration of shortfall (Proposition~\ref{prop:bounded}).

Second, under the reduced fixed-demand assumptions ($q=q_{\max}$, $D=0$, technology boost off, and no restoration), the equality balance relation declines steadily with $A_c$ until it meets the area floor. It does not fold back. The absence of a fold or critical slowing applies to this equality relation, not to every trajectory or to the separate surplus no-conversion family. With one-way conversion and no return flow, $A_c$ remains below its pre-overshoot level after conversion stops, although it no longer declines at the equality point (Table~\ref{tab:neg}). Third, recovery requires an explicit return flow; its effect depends on whether the restoration gate opens during a shortfall or a surplus.

SI~\S S7 gives a separate one-stock comparator. Its state $A$ has comparator-specific units; it is not $A_c$. With constant parameters and no debt, its equilibrium relation is $B(A)=(b+b_G\rho)A-(b_G\rho/A_{\max})A^2$, with $P=B(A)/e$. At the illustrated settings, $b_G\rho<b$, so the relation is monotone; an interior maximum would require other parameters. Its neutral equilibrium family and selected delay trajectories belong only to that formulation. They establish neither a fold nor a universal critical delay or collapse mechanism in the two-land model.

These results also clarify when to expect common early-warning indicators. Rising autocorrelation and variance are associated with approaching a fold \citep{Scheffer2001}. This reduced model reaches a minimum-area boundary instead, so those indicators need not rise first. Their absence would not establish safety. The relevant model signal is the path of $A_c$, which the aggregate alone cannot reveal.

\begin{table}[!htbp]
\centering\footnotesize
\caption{Equilibrium, trajectory, and tested delay-sensitivity properties under the stated assumptions.}
\label{tab:neg}
\begin{tabularx}{\textwidth}{>{\raggedright\arraybackslash}p{0.34\textwidth} l >{\raggedright\arraybackslash}X}
\toprule
Property & Finding & Scope \\
\midrule
Fold on the reduced frozen-demand branch & Absent & $B_{eq}(A_c)$ is strictly monotone; the balance point collides with the typed floor at $\Eceil\approx1.138\,\mathrm{gha}$ (Proposition~\ref{prop:ceiling}) \\
Critical slowing as $E\to\Eceil$ & Absent & The slow eigenvalue remains negative and constant along the reduced branch \\
Hysteresis of the reduced equality relation & Absent on this relation & There is at most one feasible equality point; fixed-demand surplus no-conversion allocations form a separate family, and the deficit-side eigenvalues are one-sided \\
Path dependence of realized $A_c$ & Present, conditional & With $R_{rc}=0$, one-way conversion leaves $A_c$ below its pre-overshoot level after shortfall closure, while $\dot A_c=0$ at equilibrium \\
Discontinuous response to endogenous demand & Not observed & The sampled minimum $A_c$ varies monotonically across tested initial populations; proximity thresholds depend on tolerance. The finite runs do not establish a general continuity theorem \\
Regeneration-delay comparison ($\tau_g$, with $\tau_p=25$\,yr) & Degenerate at 20/40\,yr & Baseline quality reaches its imposed ceiling before either delay elapses, so those paths coincide; $\tau_g=5$\,yr is distinct. This is not general delay robustness (SI Table S2) \\
\bottomrule
\end{tabularx}
\end{table}
\clearpage

\subsection{Measurement implications}

Biocapacity in gha is a value-weighted total, not an inventory of land classes. A $20\%$ increase cannot distinguish modest gains across all components from a large gain in one that offsets losses elsewhere. Proposition~\ref{prop:mask}(iii) formalizes the limit: even with fixed per-hectare yields, the total and the conserved-area constraint leave multiple feasible allocations among fast, capital, and reserve land.

A related transition-assessment study defines a linear compensatory operator and compares perfect substitutability with less-substitutable CES alternatives \citep[Sections~2 and 4 of][]{Abaee2026b}. ECOMOD's $B$ is also a linear sum, but the companion analyzes a different, finite-plan setting. There, the compensatory protocol accepts states even when no single available plan meets every component floor (Theorem~2 in \citet{Abaee2026b}). Theorem~4(iii) in \citet{Abaee2026b} identifies an exact interval of diagonal-gap states on its witness-defined region that each positive elasticity falsely certifies as safe; thus no positive-elasticity member is uniformly safe on that region. If feasibility is required throughout a transition, falling below a land-class minimum remains a violation even if the endpoint later recovers.

For the NFA, these findings concern account composition, not physical stocks. In Edition 2026, total world biocapacity account value $B$ rises $24.86\%$ from 1961 to 2023, while $B/P$ falls $52.71\%$. The Global Ecological Footprint of Production per person ($E_{\mathrm{prod}}/P$) rises $16.73\%$, and $E_{\mathrm{prod}}/B$ first exceeds one in 1973. $B$ and $E_{\mathrm{prod}}$ are annual account amounts in gha, with the annual basis implicit; per-person values are in gha per person, and $E_{\mathrm{prod}}/B$ is dimensionless.

Cropland contributes $+24.06$ percentage points to the total change; Forest biocapacity contributes $-5.95$, grazing land $-0.27$, fishing grounds $-0.42$, and built-up land $+7.44$. The non-cropland sum is nearly flat because these account movements offset one another, including the large built-up-land increase. SI~\S S5 reports the same-edition 1961--2022 comparison. None of these values measures forest biomass or condition, soil carbon, wetland extent, fish biomass, or another physical stock. In particular, a decline in Forest biocapacity does not establish a decline in forest biomass. Built-up land is an NFA accounting class, not ecological capital. The ratios, contributions, vintage, and endpoint sensitivity describe account composition only \citep{Lo2026NFA,Borucke2013,Blomqvist2013,GiampietroSaltelli2014,Galli2016,vdBerghGrazi2015}.

The model result is distinct from the NFA trends. In the frozen-demand baseline (Fig.~\ref{fig:model}a), $A_c$ falls from $1.500$ to $0.807659$\,ha ($46.16\%$), while model $B$ rises from $0.594$ to $0.825$\,gha ($39\%$). Quality $q$ rises from $0.9$ to $1.0$, moving farther above $q^*=q_{\max}/2$, where $\Ups(q)$ peaks. Its per-hectare value therefore falls from $0.05576$ to $0.05000$\,gha\,ha$^{-1}$ ($10.33\%$), and the modeled contribution $\Ups(q)A_c$ falls from $0.08364$ to $0.04038$\,gha ($51.72\%$).

With restoration off, one-way conversion leaves $A_c$ below its starting value after the shortfall closes. It then remains at the lower balance level rather than declining indefinitely. This is a conditional model result, not an observation or estimate of global physical-stock change.

We assembled no common-boundary global physical-stock series for these components and therefore report no physical-stock composite or global decline finding. The SI scoreboard is a specification—not an inventory—with candidate stock, extent, condition, and annual-flow measures and the requirements for comparing them; it contains no endpoint values. Until spatial and temporal boundaries, units, uncertainty, overlap, valuation or normalization, and permitted substitution are specified, these unlike measures should remain a nonadditive vector. SEEA Ecosystem Accounting also reports ecosystem extent and condition separately from ecosystem services and monetary presentations \citep{UN2021SEEA_EA}.

The external cropland-area and cereal-yield series provide an illustrative proxy, not direct measurements of the model's land classes. The proxy is anchored to the 1961 NFA cropland share; cereal yield is a productivity sentinel, not an all-crops measure, and the area and account boundaries do not match by construction. External cropland area $A_f$ is in hectares and index-scaled yield $\hat b_f$ in gha\,ha$^{-1}$, so $X=\hat b_fA_f$ and $C=B-X$ are account amounts in gha. The residual can absorb boundary and calibration differences, other account components, and conversion-factor changes. It is neither $\Ups(q)A_c$ nor a physical-stock measure. Physical interpretation would require independent, boundary-consistent measures of area, condition, yield, and regeneration, together with a validated mapping and valuation. Forest area could inform one part of the effective capital class, but cannot represent the entire class or its condition.

\subsection{Policy controls and limits}

We compare three controls. The rate law in Eq.~(5) already preserves the capital-land floor from safe initial states: for $y=A_c-\Acmin\ge0$, $0\le u_c\le y/\tconv$ and $R_{rc}\ge0$ imply $\dot y\ge-y/\tconv$, hence $y(t)\ge y(0)e^{-t/\tconv}$. The floor is approached asymptotically; a stop is an operational safeguard, not an additional continuous-time guarantee in this model. A ceiling on fast-land area $A_f$ does not by itself enforce the capital-land constraint. Total $B$ cannot reconstruct $A_c$ or implement arbitrary $A_c$-dependent policies, but this information limitation is not a universal control impossibility: a constant no-conversion policy also preserves the floor. Any claim of impossibility while meeting a provision objective must specify the admissible controller class, state uncertainty, and objective.

The restoration gates also behave differently. Surplus-gated restoration is inactive during a shortfall; deficit-gated restoration acts during the shortfall and can reduce provision further. This comparison echoes the distinction between protective and mobilizing institutional feedback in related work \citep{Abaee2026a}. With a common reserved share, $\sigma_f=\sigma_c=\sigma<1$, conversion is inactive while $E\le\sigma B$. With separate shares, the condition is $E\le\sigma_fY_f+\sigma_cY_c$. This reservation applies to annual human-available flow, not to spatially protected area. Half-Earth, by contrast, is a spatial conservation target \citep{Wilson2016}. The connection is only the abstract idea of reservation: the model neither maps protected areas nor evaluates that proposal. Robustness to disturbances and parameter uncertainty remains untested.

\subsection{Relation to prior work}

The ledger companion distinguishes weak and strong sustainability as operating regimes according to whether material loops close at throughput. In its terms, a positive weighted sum ranks outcomes under weak sustainability, whereas componentwise adequacy is the strong-sustainability predicate \citep[Sections~1.1 and 10.1]{Abaee2026c}. ECOMOD's $B$ is an annual biocapacity-account total. The link is an analogy about compensation, not an identity with the companion's material-flow ledger or a complete sustainability classification. The assessment companion studies aggregate measures and typed transition criteria in a separate finite-plan setting \citep{Abaee2026b}; the governance companion studies institutional feedback \citep{Abaee2026a}. The delays here represent ecological regeneration and demographic response, not institutional review.

\subsection{Limitations}

Several limitations set the scope of the results. Ecological-capital land combines unlike assets, so its single quality state, projected regeneration rule, composition diagnostic, and demand ceiling are effective model quantities—not estimates for forests, soils, or wetlands. The quality index is uncalibrated; $q_{\max}=1$ is an imposed admissibility bound, and $\rho_c$ and $\tau_g$ are scenario parameters rather than estimates from recovery studies. Most other parameters are representative rather than fitted. A cross-section cannot identify $b_c$, $b_{G,c}$, and $\rho_c$ separately. Recovery records from different systems and external measurement frameworks are precedents for future ecosystem-specific measurement and modeling only; they do not validate the present equation or parameter values. The records describe regeneration responses, not the delay $\tau_g$. We hold the per-person footprint constant to isolate feedbacks among land, biocapacity, and demand. A changing footprint would require another model extension.

The composition pattern depends on demand and on the mechanisms included here; under the stated assumptions, it is bounded. The exact shortfall switch makes stability at $S=0$ one-sided. Drawdown thresholds are reported as proximity bands because they depend on tolerance and time step. We fix debt repayment at $\eta=0.05$\,yr$^{-1}$; the zero-repayment limit ($\eta\to0$), which removes the debt equilibrium, is not analyzed. NFA accounts combine area and yield and therefore cannot identify them separately. Finally, the current-quality comparator $G_c(q(t))$ is retained only as a sensitivity check; it is not the selected projected delayed-quality model. In the displayed finite comparison it lengthens the demonstration window by $1.3$\,yr and shifts the initial-population threshold by at most $0.001$\,cap (SI~\S S3).

\section{Conclusions}\label{sec:conclusions}

Aggregate biocapacity cannot reveal how its underlying land classes are changing. In the two-land model, conversion raises total $B$ at the margin when fast-land contribution per hectare exceeds capital land's effective value ($b_f>\Ups(q)$), while $A_c$ falls if restoration is off (Proposition~\ref{prop:mask}). Other model terms determine whether total $B$ rises overall. This is a conditional accounting mechanism, not an empirical census of global physical stocks.

The reduced fixed-demand equality relation reaches the specified minimum-area boundary without a fold or critical slowing. This is a statement about that relation, not the separate fixed-demand surplus no-conversion family or every trajectory of the delayed model. Realized paths remain history-dependent: when shortfalls subside and restoration is off, $A_c$ remains below its pre-overshoot level, but it does not continue to decline at the balance point. In the representative baseline run, total model biocapacity rises $39\%$ while ecological-capital area falls $46.16\%$, its per-hectare quality term $\Ups(q)$ falls $10.33\%$, and its modeled contribution $\Ups(q)A_c$ falls $51.72\%$.

Separately, NFA Edition 2026 world account values rise $24.86\%$ between 1961 and 2023, while biocapacity per person $B/P$ falls $52.71\%$ and the production footprint per person rises $16.73\%$. Cropland contributes $+24.06$ percentage points, Forest biocapacity $-5.95$, and Built-up land $+7.44$ as an accounting class. These normalized annual account values are not observations of physical stocks or model $A_c$.

\begin{samepage}
No common-boundary global physical-stock series was assembled, so we report no composite or global physical-stock decline finding. The SI scoreboard sets out candidate measures and comparison requirements; it is not an empirical inventory. Physical measures with different units and boundaries cannot be combined into one stock total without explicit treatment of uncertainty, overlap, valuation or normalization, and allowed substitution. The auxiliary proxy residual $C=B-\hat b_fA_f$ is neither a physical-capital measure nor the model contribution $\Ups(q)A_c$; a future comparison would require independent observations and a validated mapping.\par
\end{samepage}

The practical lesson is to report component accounts alongside the aggregate. Returning to the orchard metaphor, a ledger of annual harvest alone cannot tell whether the trees, roots, and soil that support future harvests have changed. That physical record requires independent, comparable observations; it cannot be inferred from the account total. A harvest total is not a census of the orchard.

\section*{Acknowledgments}

The author thanks Global Footprint Network for making its data publicly available; Mathis Wackernagel for his thoughtful review; the editors of \emph{Ecological Modelling} for their trust and support; and Arena.ai for its ingenuity and online tools.

\section*{Declarations}

\paragraph{Data availability} All empirical inputs are publicly available from the sources cited in the manuscript. Derived analysis data and model code will be deposited in a public GitHub repository and archived in Zenodo.

\paragraph{Funding} No funding was received.

\paragraph{Competing interests} The author declares no competing interests.

\paragraph{CRediT authorship contribution statement} \textbf{Amin Abaee:} Conceptualization; Methodology; Software; Formal analysis; Writing -- original draft; Writing -- review \& editing.

\paragraph{Declaration of generative AI in scientific writing} DeepSeek AI and Claude (Anthropic) assisted with computational-reproducibility checks under the author's supervision; the author reviewed and edited the content and takes full responsibility for the publication.
